% Part: applied-modal-logics
% Chapter: temporal-logic
% Section: introduction

\documentclass[../../../include/open-logic-section]{subfiles}

\begin{document}

\olfileid{aml}{tl}{int}

\olsection{引言}

时态逻辑研究关于将要发生或已经发生之事的断言。Arthur Prior（阿瑟·普赖尔）被公认为时态逻辑的创始人，他将其称为\emph{时态逻辑}。我们在本书中对时态逻辑的阐述，将主要遵循 Prior 最初的模态进路，即将时态算子引入命题逻辑的基本框架，而命题逻辑通常将断言视为不带有时间性的。

例如，在命题逻辑中，我们可以讨论一只名叫 Beezie 的狗，它有时坐着，有时不坐着，正如狗惯常所做的那样。在经典逻辑中，断言 Beezie 正在坐着并且 Beezie 不在坐着将是矛盾的。但显然两者都可能为真，只是不在同一时刻为真；为语言增加时态算子可以让我们相对容易地表达这一断言。时态算子的引入，也让我们能够解释某些推理的有效性，例如从“Beezie 将得到一个零食或一个球”推出“Beezie 将得到一个零食，或者 Beezie 将得到一个球”。

然而，时态逻辑引发了大量哲学问题，可能促使我们倾向于采用某一种而非另一种时态逻辑框架。例如，一个未来偶然语句是关于未来的、既非必然也非不可能的语句。如果我们说“Richard 明天会去杂货店”，我们就是在表达一个关于尚未发生之事的断言，其真值是有争议的。事实上，这个断言根本上能否被\emph{指派}一个真值，本身亦存争议。如果我们是严格的决定论者，那么我们或许可以接受这个语句实际上是有真值或假值的，甚至在该事件预定发生之前——仅仅是我们还不知道它的真值而已。相反，我们可能相信一个真正开放的未来，其中未来偶然语句的真值是未定的。

事实证明，许多关于时间结构和本质的承诺，都内置于我们对时态逻辑模型和框架的选择之中。例如，我们可能会问自己，应当构建时间是线性的、分支的甚至是循环的模型。我们可能需要决定，我们的时态模型是否要有起点和终点，以及时间是应该用离散的瞬间来表示，还是作为一个连续统来表示。

\end{document}